\chapter{The complete causal statement}

The most defensible formulation is:

\begin{quote}
Divergent natural selection can change both the curvature of local fitness surfaces and the distribution of standing genetic variance. These changes rotate the leading invariant subspaces of the symmetric response operator \(\Sop=\G^{1/2}\Hh\G^{1/2}\) and, when \(\G\) evolves, also deform the map from whitened coordinates back to phenotype space. Under smooth deterministic divergence, common-metric squared channel distance increases quadratically at early time, with an exact coefficient \(\normF{\mathbf K_M}^2\) determined by cross-channel forcing, spectral gaps, and the changing square-root maps. Channel distance becomes reproductive isolation only when parental covariance structures are coadapted to their own fitness curvatures, so that recombinant hybrids express variance in directions strongly penalised in parental habitats. This early ecological barrier can reduce effective gene flow and thereby facilitate the later combinatorial accumulation of intrinsic Dobzhansky--Muller incompatibilities.
\end{quote}

The theory is therefore neither a replacement for Lande's response equation nor a replacement for Orr's snowball. It is a mathematical bridge between them.

\begin{figure}[H]
\centering
\includegraphics[width=0.96\textwidth]{figures/figure_geometric_snowball_concept.pdf}
\caption{Integrated conceptual view. Ecological divergence rotates parental channels; the resulting overlap deficit begins slowly and then accelerates. The fitness interpretation requires the covariance--curvature calibration derived in this document.}
\label{fig:integratedConcept}
\end{figure}

